\documentclass[a4paper,11pt]{article}

\usepackage[T1]{fontenc}
\usepackage{lmodern}
\usepackage[utf8]{inputenc}
\usepackage[ngerman]{babel}
\usepackage{csquotes}
\usepackage[backend=biber,style=authoryear]{biblatex}
\usepackage{parskip}
\usepackage{microtype}
\usepackage{booktabs}
\usepackage{tabularx}
\addbibresource{literatur.bib}

\begin{document}

% ---------------------------------------------


\section{Unit-Tests für Lexer, Parser und Codegenerator}

Unit-Tests überprüfen gezielt einzelne Einheiten, wie Funktionen und Komponenten einer Software \parencite{meszaros2007}. Sie sind besonders hilfreich, wenn sie auf eine bestimmte Phase des Compilationsprozesses abzielen, da dadurch präziser die Ursache eines Fehlers identifiziert werden kann.  

\subsection{Lexer}

\begin{verbatim}
EXPECT_TRUE(Lex("#include <\\\\\n").empty());
\end{verbatim}

\subsection{Parser}


\subsection{Codegenerator}


\section{Differential Testing}

Ein Programm sollte bei gleichem Quellcode auch mit unterschiedlichen Compilern
und Compileroptionen das gleiche beobachtbare Verhalten zeigen. Voraussetzung
ist, dass der Quellcode kein undefiniertes Verhalten enthält und relevante
Unterschiede zwischen den Zielplattformen berücksichtigt werden.

Dabei werden folgende Varianten unterschieden:

\begin{itemize}
  \item Cross-Compiler Testing \parencite{terres2025}
  \item Cross-Version Testing \parencite{terres2025}
  \item Cross-Optimization Testing \parencite{terres2025}
  \item Cross-Architecture Testing \parencite{elnaggar2025}
\end{itemize}

\subsection{Cross-Compiler Testing}

Unterschiedliche Compiler, wie beispielsweise GCC und Clang, arbeiten intern
unterschiedlich, müssen aber die Semantik eines Programms bei der Übersetzung
erhalten.

\subsection{Cross-Version Testing}

Verschiedene Versionen eines Compilers werden miteinander verglichen, um
Regressionen aufzudecken. Bei unveränderten Sprachregeln und vergleichbaren
Compileroptionen sollte das beobachtbare Programmverhalten erhalten bleiben.

\subsection{Cross-Optimization Testing}

Compiler bieten unterschiedliche Optimierungsstufen an, beispielsweise
\texttt{-O1}, \texttt{-O2} und \texttt{-O3}. Dabei können unter anderem
Kontrollstrukturen verändert und Ausdrücke vereinfacht werden. Das definierte
beobachtbare Programmverhalten muss dabei erhalten bleiben.

\subsection{Cross-Architecture Testing}

Das Programmverhalten wird auf unterschiedlichen Zielarchitekturen verglichen,
beispielsweise x86, ARM, MIPS und RISC-V. Dabei müssen plattformabhängige
Eigenschaften wie die Größe von Datentypen und die Byte-Reihenfolge
berücksichtigt werden.

Abweichende Ergebnisse können auf Compilerfehler hinweisen.
Der Vergleich allein zeigt jedoch weder eindeutig, \emph{welche}
Konfiguration fehlerhaft ist, noch, \emph{warum} der Fehler auftritt.

Csmith \parencite{csmith} generiert zufällige Programme in C. Diese können mit
verschiedenen Compilern und Compileroptionen getestet werden. Die
Ergebnisse werden anschließend verglichen, um mögliche Fehler aufzudecken.

\begin{verbatim}
csmith > random2.c
gcc random2.c -I$HOME/csmith/include -o random2_gcc
clang random2.c -I$HOME/csmith/include -o random2_clang
./random2_gcc > gcc_output.txt
./random2_clang > clang_output.txt
diff -u gcc_output.txt clang_output.txt
\end{verbatim}

Der obige Ausschnitt zeigt Befehle für eine Unix-Shell. Ein zufälliges
C-Programm wird generiert, mit GCC und Clang kompiliert und anschließend
ausgeführt. Die Ausgaben werden in Textdateien gespeichert und mit
\texttt{diff} verglichen. Der Include-Pfad muss zur lokalen
Csmith-Installation passen.

Die erzeugten Binärdateien dürfen sich unterscheiden, da die Compiler intern
unterschiedlich arbeiten. Unter vergleichbaren Bedingungen muss das definierte
beobachtbare Verhalten und damit auch die verglichene Ausgabe übereinstimmen.

\section{Gute Compiler-Fehlermeldungen}

Hilfreiche Compilerfehlermeldungen erleichtern es, einen Fehler zu lokalisieren
und zu beheben. Die folgende Tabelle enthält illustrative Beispiele; der genaue
Wortlaut hängt unter anderem von der Compilerversion und dem Quellcode ab.

\begin{table}[htbp]
  \centering
  \caption{Illustrative Beispiele für Compilerfehlermeldungen}
  \label{tab:compiler-fehlermeldungen}
  \small
  \renewcommand{\arraystretch}{1.2}
  \begin{tabularx}{\linewidth}{@{}l>{\raggedright\arraybackslash}Xl@{}}
    \toprule
    \textbf{Angabe} & \textbf{Beispiel} \\
    \midrule
    Schweregrad & \texttt{Error/Warn/Info/Critical}  \\
    Fehlerart & \texttt{Syntax Error} \\
    Kontext & \texttt{In der Funktion main()} \\
    Betroffene Position & \texttt{main.py, Zeile 10, Spalte 3} \\
    Betroffener Quelltext & \texttt{print 1} \\
    Vorschlag zur Behebung & \texttt{Füge Klammern hinzu} \\
    \bottomrule
  \end{tabularx}
\end{table}

Um sicherzustellen, dass Compilerfehlermeldungen stets konsistent und hilfreich sind, können auch diese mit Unit-Tests überprüft werden. 




% ---------------------------------------------

\clearpage
\printbibliography

\end{document}
